\section*{Hoofdstuk \ref{ch:b3:behavioural-ecology} --- Gedragsecologie}

\begin{solution}{exo:b3:behavioural-ecology:1}
Mechanisme: de danseres zet de hoek tussen voedsel en zon, tijdens de
vlucht gemeten, om in de hoek met de zwaartekracht op de raat, en
codeert de afstand uit de optische stroom van de reis als de duur van
de loop. Ontwikkeling: grotendeels aangeboren --- bijen die zijn
opgekweekt zonder ooit een dans te zien dansen correct, met kleine
geleerde bijstellingen. Functie: zij rekruteert nestgenoten naar een
lonende bron en verhoogt zo de opbrengst van het volk. Geschiedenis:
afgeleid van geritualiseerde aanzetten tot vertrek naar het voedsel;
dwerghoningbijen dansen nog op een horizontaal vlak en wijzen
rechtstreeks naar de bron, de voorouderlijke vorm.
\end{solution}

\begin{solution}{exo:b3:behavioural-ecology:2}
Een \omterm{def:b3:behavioural-ecology:tinbergen}{vast handelingspatroon} is een gestandaardiseerde reeks die door een
eenvoudige aanwijzing, de \omterm{def:b3:behavioural-ecology:tinbergen}{sleutelprikkel}, wordt vrijgemaakt en die
eenmaal begonnen wordt afgemaakt: het eirollen van de gans, vrijgemaakt
door een eivormig voorwerp; het pikken van het meeuwenkuiken,
vrijgemaakt door een rode vlek op een snavelachtige vorm. Het stokje
met drie rode banden, waarnaar vaker werd gepikt dan naar een echte
kop, liet zien dat het vrijmakende mechanisme een filter is dat op een
kenmerk is afgestemd --- rood contrast, langgerektheid --- en geen
herkenning van de ouder, en dat het kenmerk overdrijven het antwoord
overdrijft.
\end{solution}

\begin{solution}{exo:b3:behavioural-ecology:3}
Verlaat een plek wanneer het tempo van winst erin is gedaald tot het
gemiddelde tempo voor het hele leefgebied, de reis inbegrepen. De
reistijd verdubbelen verlaagt dat gemiddelde, dus blijft de
foerageerder in elke plek langer, put hij elke plek grondiger uit, en
eindigt hij met een lager algemeen tempo.
\end{solution}

\begin{solution}{exo:b3:behavioural-ecology:4}
Een ESS is een strategie die, zodra vrijwel iedereen haar speelt, door
geen enkel zeldzaam alternatief kan worden overtroffen. Louter Duiven
is er geen: één enkele Havik wint elke wedstrijd en verdient $V$ tegen
de $V/2$ van de Duiven. In het mengsel verdient elke strategie
$V(1 - V/C)/2$, minder dan de $V/2$ die louter Duiven zou geven, omdat
een deel van de wedstrijden gevechten zijn die verwonden; niemand kan
er in zijn eentje op verbeteren, dus zit iedereen met het verlies.
\end{solution}

\begin{solution}{exo:b3:behavioural-ecology:5}
$V = 6$, $C = 10$: $p^{*} = 0.6$; opbrengst $6(1 - 0.6)/2 = 1.2$ voor
elk, tegen $3$ bij louter Duiven. $V = 12 > C$: Havik is de zuivere
ESS, opbrengst $(12 - 10)/2 = 1$, tegen $6$ bij louter Duiven --- de
wedstrijd vernietigt vijf zesde van de waarde.
\end{solution}

\begin{solution}{exo:b3:behavioural-ecology:6}
$\tau = T_{0}$: $t^{*} = 1.15\times 2 = \qty{2.3}{min}$, deel $1 -
\mathrm{e}^{-1.15} = 0.68$, $R = 0.68A/4.3 = 0.16A$ per minuut. $\tau =
4T_{0}$: $t^{*} = 2.0\times 2 = \qty{4}{min}$ (nauwkeuriger $3.9$),
deel $0.86$, $R = 0.86A/12 = 0.072A$ per minuut.
\end{solution}

\begin{solution}{exo:b3:behavioural-ecology:7}
Halfbroers en halfzussen $1/4$; grootouder en kleinkind $1/4$; volle
neven en nichten $1/8$. Een haplodiploïde werkster en haar broer: hij
is haploïd en heeft alleen de genen van hun moeder; van de genen van de
werkster kwam de helft van de moeder, en de helft daarvan zijn de
kopieën die hij ontving, dus $r = 1/4$ vanuit haar gezichtspunt (vanuit
het zijne is het $1/2$: de verwantschap is niet symmetrisch wanneer de
ploïdie verschilt). Koningin tot zoon: al zijn genen zijn de hare, maar
slechts de helft van de hare zit in hem: $r = 1/2$ van haar kant, $1$
van de zijne.
\end{solution}

\begin{solution}{exo:b3:behavioural-ecology:8}
$n\,r\,b > c$ met $b = 0.01$ en $c = 0.02$: $n\,r > 2$. Volle broers en
zussen, $r
= 1/2$: $n \ge 5$. Nichtjes, $r = 1/4$: $n \ge 9$. Niet-verwant: nooit
--- en daarom roepen wegtrekkende mannetjes niet.
\end{solution}

\begin{solution}{exo:b3:behavioural-ecology:9}
De paringen verdubbelen per \qty{25}{cm}, de overleving daalt een
vijfde: $L = 25$: $0.5\times 1.2 = 0.6$; $50$: $1$; $75$: $2\times 0.8 = 1.6$;
$100$: $4\times 0.6 = 2.4$; $125$: $8\times 0.4 = 3.2$; $150$: $16\times 0.2 =
3.2$. Op deze voorwaarden blijft het product ver boven
de natuurlijke \qty{50}{cm} stijgen, dus mist het model kosten: de
overlevingsstraf van een lange staart in de vlucht is steiler dan
lineair, het antwoord van de vrouwtjes verzadigt, het laten groeien van
de veren kost energie, en de verdubbeling van het aantal nesten bij
Andersson werd over één seizoen gemeten en niet over een leven. De
natuurlijke staart ligt waar de werkelijke kosten bijten.
\end{solution}

\begin{solution}{exo:b3:behavioural-ecology:10}
Vergelder: tegen Havik $(V - C)/2$, tegen Duif $V/2$, tegen Vergelder
$V/2$. In een populatie van louter Vergelders treft een zeldzame Havik
Vergelders aan die terugvechten, en verdient hij $(V - C)/2 < V/2$: hij
kan niet binnendringen wanneer $V < C$. Een zeldzame Duif verdient
$V/2$, precies de opbrengst van de Vergelders, dus zijn Duiven neutraal
en kunnen zij naar binnen drijven; eenmaal algemeen laten zij Haviken
toe (een Havik verdient $V$ tegen een Duif), en heeft het spel met drie
strategieën geen eenvoudige ESS --- een eerste voorbeeld van hoe het
toevoegen van een optie een stabiel mengsel kan verstoren.
\end{solution}

\begin{solution}{exo:b3:behavioural-ecology:11}
Hield iedereen het precies $t_{0}$ vol, dan zou een mutant die $t_{0} +
\varepsilon$ volhoudt elke wedstrijd winnen voor verwaarloosbaar veel
extra kosten, en zou een populatie van zulke mutanten door $t_{0} + 2\varepsilon$
worden verslagen, en zo voort: geen vaste tijd is stabiel, dus moet de
ESS een willekeurige tijd zijn. Bij een exponentiële verdeling is de
kans om nog $s$ langer vol te houden, gegeven dat het al $t$ duurt,
$\mathrm{e}^{-s/m}$ wat $t$ ook is --- de geheugenloosheid --- zodat
een tegenstander die nog vertoont niets prijsgeeft over hoeveel langer
hij doorgaat, en geen enkele strategie de verstreken tijd kan uitbuiten.
\end{solution}

\begin{solution}{exo:b3:behavioural-ecology:12}
Werksters waarderen een zus op $3/4$ en een broer op $1/4$, dus zou de
kolonie onder het bewind van de werksters drie maal zoveel in vrouwtjes
als in mannetjes moeten investeren, een verhouding $3{:}1$ naar
investering; de koningin, aan beide even verwant, geeft de voorkeur aan
$1{:}1$. Trivers en Hare (1976) vonden investeringsverhoudingen rond
$3{:}1$ bij mieren met één koningin die één maal had gepaard: de
werksters winnen. Meervoudig paren verlaagt de gemiddelde verwantschap
van een werkster met haar zussen naar $1/4$ tot $1/2$ en duwt het
optimum van de werksters naar $1{:}1$; kolonies van soorten met
veelvuldig gepaarde koninginnen investeren inderdaad gelijkmatiger,
zoals het betoog voorspelt.
\end{solution}

\begin{solution}{pb:b3:behavioural-ecology:1}
\textbf{1.} $R(t) = A(1 - \mathrm{e}^{-t/T_{0}})/(\tau + t)$; vertrek
wanneer $g'(t) = R(t)$: $(A/T_{0})\mathrm{e}^{-t/T_{0}} = A(1 - \mathrm{e}^{-t/T_{0}})/
(\tau + t)$.
\textbf{2.} Vermenigvuldig met $(\tau + t)\mathrm{e}^{t/T_{0}}/A$: $(\tau + t)/T_{0}
= \mathrm{e}^{t/T_{0}} - 1$. Met $\tau = T_{0}$ en $x = t/T_{0}$: $2 + x =
\mathrm{e}^{x}$; $x = 1.15$ geeft $3.15$ tegen $3.16$.
\textbf{3.} $\tau = 4T_{0}$: $5 + x = \mathrm{e}^{x}$; $x = 1.94$ geeft
$6.94$ tegen $6.96$; $t^{*} \approx 1.94\,T_{0}$, zeg $2.0$.
\textbf{4.} Dicht: $t^{*} = \qty{3.5}{min}$, $60(1 - \mathrm{e}^{-1.15}) =
\qty{41}{kJ}$, \qty{68}{\%}, $R = 41/6.5 = \qty{6.4}{kJ/min}$. IJl:
$t^{*} = \qty{5.8}{min}$, \qty{51}{kJ}, \qty{86}{\%}, $R = 51/17.8 =
\qty{2.9}{kJ/min}$.
\textbf{5.} Vijf minuten: $g(5) = 60(1 - \mathrm{e}^{-5/3}) = \qty{49}{kJ}$.
Dicht: $49/8 = \qty{6.1}{kJ/min}$ (\qty{96}{\%} van het optimum); ijl:
$49/17 = \qty{2.9}{kJ/min}$ (\qty{99}{\%}). Een grove regel verliest een
paar procent --- goed genoeg opdat de selectie er genoegen mee heeft
genomen.
\textbf{6.} De tussentijd tussen vangsten wordt langer naarmate de plek
uitgeput raakt, dus is een vaste opgeeftijd een vast marginaal
vangsttempo --- het criterium van de stelling, dat in elke plek
hetzelfde is. Zij faalt omdat de juiste drempel van het leefgebied
afhangt: met een vaste \qty{20}{s} vertrekt het dier te vroeg waar de
reis lang is en te laat waar zij kort is, tenzij de opgeeftijd zelf aan
de recente ervaring wordt aangepast --- wat foerageerders doen.
\textbf{7.} Havik tegen Havik $(10 - 30)/2 = -10$, Havik tegen Duif
$10$, Duif tegen Havik $0$, Duif tegen Duif $5$; $p^{*} = V/C = 1/3$.
\textbf{8.} $W_{H} = 10 - 20p$, $W_{D} = 5(1 - p)$. $p = 0.1$: $8$ tegen
$4.5$, Haviken breiden uit. $p = 1/3$: elk $3.33$, evenwicht. $p = 0.6$:
$-2$ tegen $2$, Duiven breiden uit.
\textbf{9.} $3.33$ bij de ESS tegen $5$ bij louter Duiven: een derde van
de waarde van elke wedstrijd gaat verloren aan de mogelijkheid van
vechten.
\textbf{10.} $V = 40 > C$: Havik is de zuivere ESS; elke wedstrijd is
een gevecht. In zware jaren zouden er meer gevechten en meer
verwondingen moeten zijn --- en die zijn er ook.
\textbf{11.} Een Havik die een Bourgeois ontmoet treft die de helft van
de tijd als bewoner aan (een gevecht, $(V - C)/2$) en de helft van de
tijd als indringer (een walk-over, $V$): gemiddeld $(3V - C)/4$.
Bourgeois verdient $V/2$ onder zijns gelijken. Havik dringt alleen
binnen als $(3V - C)/4 > V/2$, d.w.z.\ $V > C$. Het ``de bewoner wint''
van de vlinder is Bourgeois: een goedkope afspraak, stabiel omdat
vechten meer kost dan een zonnevlek waard is --- en wanneer beide
zichzelf voor bewoner houden, spelen beide Havik.
\textbf{12.} Wat het beste is om te doen hangt af van wat alle anderen
doen, dus komt de populatie uit waar de strategieën even goed lonen en
niet bij het beste van enig individu.
\textbf{13.} $n\bar r\,b > c$: $n\bar r > c/b = 3$.
\textbf{14.} Acht volle broers en zussen: $n\bar r = 4 > 3$, de dienst
loont. Acht vreemden: $0$, nooit. Acht halfbroers en halfzussen: $2 < 3$,
nee.
\textbf{15.} $c = 0.01$: de drempel wordt $n\bar r > 1$ --- nu volstaan
halfbroers en halfzussen en zelfs een paar volle broers en zussen. Het
dier met de kleinste kosten is het dier voor wie de regel is voldaan,
dus staan de goed gevoede dieren op wacht terwijl de hongerige graven.
\textbf{16.} $k$ zussen zijn $\tfrac{3}{4}k$ waard, $k$ dochters
$\tfrac{1}{2}k$: verkies zussen, met $3{:}2$.
\textbf{17.} Een willekeurige zus is een volle zus ($3/4$) met kans
$1/3$ en een halfzus ($1/4$) met kans $2/3$: $\bar r = 0.25 +
0.17 = 0.42 < 0.5$. De voorkeur keert om: dochters verslaan nu zussen,
dus kan de verwantschap alleen de steriele werksters bij polyandrische
soorten niet verklaren --- daar moeten de ecologie en het toezicht op
het niveau van de kolonie bij.
\textbf{18.} De regel vergt alleen dat $r$ onder de geholpen dieren
gemiddeld hoger is dan in de populatie; elke aanwijzing die met
verwantschap samenhangt volstaat --- in de buurt van het geboortehol
blijven, helpen wie in hetzelfde hol is opgegroeid. De regel van een
grondeekhoorn kan zijn ``roep dicht bij huis'', wat
verwanten selecteert zonder er ook maar één te herkennen.
\textbf{19.} $W = (L/20)(1 - L^{2}/10^{4})$: $L = 20$: $0.96$; $40$:
$1.68$; $60$: $1.92$; $80$: $1.44$.
\textbf{20.} $W' = \tfrac{1}{20}(1 - 3L^{2}/10^{4}) = 0$: $L^{*} = 100/\sqrt{3}
= \qty{58}{cm}$, $W = 2.89\times\tfrac{2}{3} = 1.92$.
\textbf{21.} $L^{*} = 70/\sqrt{3} = \qty{40}{cm}$: korter. Het optimum
ligt waar de marginale winst aan paringen gelijk is aan het marginale
verlies aan overleving, dus trekt een zwaardere predatiekost de staart
in, en zouden populaties onder verschillende predatie in hun sieraad
moeten verschillen --- zoals guppy's en zwaarddragers doen.
\textbf{22.} Beide. De \omterm{def:b3:behavioural-ecology:sexual-selection}{op hol geslagen selectie} voorspelt dat de
voorkeur het kenmerk vooruitsnelt, dat bij het overlevingsoptimum
stopt; de handicap voorspelt dat vrouwtjes langer verkiezen omdat
alleen een fit mannetje dat draagt, opnieuw voorbij wat de meeste
kunnen laten groeien. Het experiment toont een voorkeur voorbij het
optimum en beslist niet tussen de verklaringen.
\textbf{23.} Was de staart gratis, dan zou elk mannetje de langste
kunnen laten groeien, zou hij niets over zijn kwaliteit zeggen, en zou
een voorkeur ervoor worden uitgebuit en wegkwijnen. Een kost die
zwaarder weegt voor een mannetje in slechte conditie maakt de lange
staart alleen voor de fitte betaalbaar, en dan kiest de voorkeur de
fitheid: de kost is wat het signaal eerlijk maakt.
\textbf{24.} Met acht maal de opbrengst per paring stijgt de fitness van
een mannetje steil met het aantal paringen en die van een vrouwtje
nauwelijks; mannetjes worden geselecteerd om om paringen te wedijveren
--- de wedstrijden Havik--Duif en hun wapens --- en vrouwtjes, voor wie
elk van de weinige nakomelingen telt, worden geselecteerd om te kiezen,
wat de sieraden van de mannetjes doet lonen. Eén scheefheid in de winst
per paring brengt zowel het vechten als de opsmuk voort.
\textbf{25.} $t^{*} \approx \qty{3.5}{min}$ in het dichte leefgebied; $p^{*}
= 1/3$; de schildwachtdienst loont wanneer $n\bar r > 3$.
\end{solution}
